% Не більше 4 сторінок. Мета — практичний результат, не «дослідження...».
\structural{ВСТУП}

\textbf{Актуальність теми.} \todo{критичний аналіз і порівняння з відомими розв'язаннями в контексті програмної інженерії} \cite{example-docs}.

\textbf{Мета проєкту.} \todo{мета як очікуваний практичний результат}.

\textbf{Об'єктом дослідження} є \todo{...}.

\textbf{Предметом дослідження} є \todo{...}.

Для досягнення поставленої мети необхідно виконати такі завдання:
\begin{itemize}
\item проаналізувати предметну область та існуючі рішення;
\item сформувати функціональні й нефункціональні вимоги;
\item спроєктувати архітектуру програмного забезпечення;
\item розробити програмне забезпечення;
\item провести тестування;
\item описати розгортання та супровід.
\end{itemize}

\textbf{Основні проєктні рішення.} \todo{...}

\textbf{Практичне значення отриманих результатів} \todo{...}.

\textbf{Можливі галузі застосування} \todo{...}.
